\hypertarget{appendix-b-audit-details}{%
\section{Audit details}\label{appendix-b-audit-details}}

\hypertarget{b.1-data-and-likelihood}{%
\subsection{Data and likelihood}\label{b.1-data-and-likelihood}}

Planck 2018 compressed distance priors, DESI DR2 compressed BAO and
Pantheon+ supernova distances, \(n=1604\) points in the joint fit. The
supernova absolute intercept is profiled analytically and the
sound-horizon convention is fixed as in the pipeline's validation. The
background is computed with a compressed scipy-based pipeline;
multi-parameter dark-energy fits use nested sampling (dynesty). No
Boltzmann solver is used anywhere in the geometry audits.

\hypertarget{b.2-pre-registration-summary}{%
\subsection{Pre-registration
summary}\label{b.2-pre-registration-summary}}

Pre-registered before the audit tables were generated, using existing
data products and baseline fits:

\begin{itemize}
\tightlist
\item
  \textbf{Models (B1).} \(\Lambda\)CDM (\(k=4\)); \(w_0w_a\)CDM
  (\(k=6\)); rewrite \(E^2(z)=\Omega_m(1+z)^3+\Omega_r(1+z)^4+A(1+z)^m\)
  with \(A\ge0\) and bounded \(m\) (\(k=6\)). Readouts \(\chi^2\), AIC,
  BIC and the posterior of \(m\). Constant-compatibility criterion:
  \(|m|\le\min(2\sigma_m,0.25)\), evaluated at the best-fit \(m\).
\item
  \textbf{Chart comparison (B1b).} The best-fit \(w_0w_a\) history as
  \(w(z)\) and as \(\Delta_H(z)\).
\item
  \textbf{Staging statistic (B2).} \(S\) as in Section 8.3 over the
  splits BAO only, supernovae only, supernovae with \(z<0.35\) and with
  \(z\ge0.35\); Planck is excluded because the audit targets supernova
  and BAO route dependence. Null: mock \(\Lambda\)CDM realizations
  through the identical pipeline. All splits reported with bound and
  convergence flags.
\item
  \textbf{Route residuals (B3).} Unsigned parameter residuals of the
  supernova-only and BAO-only fits against the joint fit.
\item
  \textbf{Quarantine.} Fitted amplitudes are in-sample parameters, never
  predictions.
\end{itemize}

\hypertarget{b.3-model-and-chart-results}{%
\subsection{Model and chart
results}\label{b.3-model-and-chart-results}}

\begin{longtable}[]{@{}lllllll@{}}
\toprule
model & \(k\) & \(\chi^2\) & AIC & BIC & \(\Delta\)AIC &
\(\Delta\)BIC\tabularnewline
\midrule
\endhead
\(\Lambda\)CDM & 4 & 1455.80 & 1463.80 & 1485.32 & 0 & 0\tabularnewline
\(w_0w_a\)CDM & 6 & 1442.17 & 1454.17 & 1486.45 & \(-9.63\) &
\(+1.13\)\tabularnewline
rewrite & 6 & 1453.66 & 1465.66 & 1497.94 & \(+1.85\) &
\(+12.61\)\tabularnewline
\bottomrule
\end{longtable}
\addtocounter{table}{-1}

\(\Lambda\)CDM: \(\Omega_m=0.303\). \(w_0w_a\)CDM: \(w_0=-0.795\),
\(w_a=-0.673\). Rewrite: best fit \(A=0.699\), \(m=0.068\), with
\(\sigma_m=0.074\); posterior median \(m=0.131\) (16th and 84th
percentiles \(0.057\), \(0.204\)). Chart comparison: crossing at
\(z_\star\approx0.4388\), the rounded analytic root for the stored
best-fit decimals (a grid interpolation gave \(0.43885\)); \(\Delta_H\)
between \(0.4985\) and \(0.7596\) on the tabulated range, finite and
smooth throughout. All 251 tabulated chart rows agree with the analytic
density to \(10^{-12}\).

\hypertarget{b.4-staging-calibration-and-amendments}{%
\subsection{Staging calibration and
amendments}\label{b.4-staging-calibration-and-amendments}}

\textbf{Original alternatives.} A rewrite \(\Delta_H=A(1+z)^m\) injected
coherently into all probes, at amplitudes up to \(A=5.6\), was never
detected. The real score also exceeded the threshold at that stage
(\(9.63\) against \(8.67\)), and was not interpreted, because the audit
had not been shown to detect anything.

\textbf{Amended alternatives (logged).} A supernova-only tilt
\(\mu\to\mu+\varepsilon(z/0.5)\) and a BAO-only tilt
\(\alpha\to\alpha(1+\varepsilon(z-0.8))\) with
\(\varepsilon\in\{0.01,0.02,0.05,0.10,0.20\}\), 20 realizations each,
and the same statistic and null. Detection counts at the rank-rule
cutoff \(8.879\):

\begin{longtable}[]{@{}llllll@{}}
\toprule
\(\varepsilon\) & 0.01 & 0.02 & 0.05 & 0.10 & 0.20\tabularnewline
\midrule
\endhead
supernova tilt & 1/20 & 0/20 & 1/20 & 2/20 & 17/20\tabularnewline
BAO tilt & 5/20 & 3/20 & 7/20 & 20/20 & 20/20\tabularnewline
\bottomrule
\end{longtable}
\addtocounter{table}{-1}

\textbf{Threshold rule.} The originally declared threshold was the
interpolated empirical 95th percentile of the 20 null scores, \(8.328\).
Interpolated percentiles can exceed their nominal false-alarm rate even
for exchangeable scores: for a uniformly random permutation of the
equally spaced scores \(0,\ldots,20\), with twenty used for calibration
and one held out, the interpolated 95th-percentile rule rejects with
probability \(2/21>1/20\). The conservative rank rule rejects only
scores above the \(\lceil(n+1)(1-\alpha)\rceil\)-th smallest null score
(never rejecting if that rank exceeds \(n\)), here the maximum,
\(8.879\); under exchangeability its rejection probability is at most
\(1/21\). Both the finite counting law and the counterexample are
mechanized. With either threshold the real score \(9.630\) fires, and
the detection-grid values remain \(0.20\) and \(0.10\).

\textbf{Split fits.}

\begin{longtable}[]{@{}llllll@{}}
\toprule
split & data points & \(A\) & \(m\) & bound flag & split
score\tabularnewline
\midrule
\endhead
joint & 19 & 0.817 & \(-0.454\) & none & ---\tabularnewline
BAO only & 13 & 0.701 & 0.053 & none & 6.35\tabularnewline
supernovae only & 6 & 0.400 & \(-6.0\) & \(m\) at lower bound &
7.10\tabularnewline
supernovae, \(z<0.35\) & 3 & 0.398 & \(-6.0\) & \(m\) at lower bound &
6.33\tabularnewline
supernovae, \(z\ge0.35\) & 3 & 8.0 & 1.530 & \(A\) at upper bound &
9.63\tabularnewline
\bottomrule
\end{longtable}
\addtocounter{table}{-1}

The reduced vector consists of thirteen BAO points and six compressed
supernova summaries, with \(\Omega_m\) fixed at \(0.3027\). Reduced
vectors are fragile for shape fits, which is why several fits reach a
bound; the \(0.05\) term in the denominator of \(S\) limits the effect
of small split variances. An early rewrite fit on this reduced vector
(\(\chi^2/{\rm dof}\approx6.7\), with bound hits) was not comparable to
the baselines; Section 8.1 reports the refit through the full
likelihood.

\hypertarget{b.5-route-residuals}{%
\subsection{Route residuals}\label{b.5-route-residuals}}

\begin{longtable}[]{@{}lll@{}}
\toprule
route pair & parameter & unsigned residual\tabularnewline
\midrule
\endhead
supernovae only vs joint & \(\Omega_m\) & 0 (fixed)\tabularnewline
supernovae only vs joint & \(A\) & 0.417\tabularnewline
BAO only vs joint & \(\Omega_m\) & 0 (fixed)\tabularnewline
BAO only vs joint & \(A\) & 0.116\tabularnewline
\bottomrule
\end{longtable}
\addtocounter{table}{-1}

\hypertarget{b.6-growth-audit}{%
\subsection{Growth audit}\label{b.6-growth-audit}}

Probes: KiDS-Legacy COSEBIs (126 points, S/N 28.61), HSC-Y3
Fourier-space cosmic shear (170 points, S/N 37.39), ACT, SPT and Planck
CMB lensing (35 bandpowers, S/N 61.35), with BOSS and DESI full-shape
clustering for the mild regime. Heterogeneity proxy
\(H=\langle[\log(P_{\rm nl}/P_{\rm lin})]_+\rangle_W\), evaluated in one
posterior-mean cosmology. Shell census (slack \(H\le0.05\) / mild /
strong \(H\ge0.3\)): KiDS \(0/3/0\), HSC \(0/0/4\), CMB lensing
\(2/1/0\); full-shape about \(H\le0.11\) at \(k\lesssim0.2\). The
within-data-set comparison computed no price because its support
condition failed. In separate HSC diagnostics with local CCL-based
linearizations over posterior draws, linear-model shell prices are
\(8.29\), \(8.80\), \(16.50\) and \(10.57\), and nonlinear models with
and without nuisance parameters give zero point prices. The nonlinear
controls use finite injections and return injected point prices of
\(3.9\). The linear controls already exceed the threshold before
injection; their injection search ends at the negligible scale
\(2^{-80}\), so they do not demonstrate incremental sensitivity to an
injected signal. The fitted geometry inputs of these diagnostics are not
independent out-of-sample data.
